\documentclass{beamer}
\usetheme{Madrid}
\usepackage{amsmath, amssymb, amsthm}
\usepackage{graphicx}
\usepackage{listings}
\usepackage{gensymb}
\usepackage[utf8]{inputenc}
\usepackage{hyperref}
\usepackage{gvv}
\usepackage{tikz}
\lstset{
  language=Python,
  basicstyle=\ttfamily\small,
  keywordstyle=\color{blue},
  stringstyle=\color{orange},
  numbers=left,
  numberstyle=\tiny\color{gray},
  breaklines=true,
  showstringspaces=false
}
\usetikzlibrary{decorations.pathmorphing}

\title{Question-9.4.35}
\author{EE25BTECH11022 - sankeerthan}
\date{}
\begin{document}

\frame{\titlepage}

\begin{frame}
\frametitle{Question}
A train travels a distance of $480km$ at a uniform speed. If the speed had been 8$\frac{km}{h}$ less, then it would have taken 3 hours more to cover the same distance. We need to find the speed of the train.
\end{frame}

\begin{frame}
\frametitle{Solution}
Let the unknown speed be x.
\begin{align}
\vec{t} =480\myvec{\frac{1}{x-8} \\\frac{1}{x}} \\
\vec{A}^T\vec{X} =3,\text{where}\vec{A} =\myvec{1 \\ -1},
\vec{X} =\myvec{\frac{480}{x-8} \\\frac{480}{x}}
\end{align}
Multiply both sides by $x(x-8)$:
\begin{align}
480x-480(x-8)&=3x(x-8) \\
480x-480x+3840 &=3x^2-24x \\
3x^2-24x-3840 &=0 \\
\myvec{x^2 & x & 1}\myvec{3 \\ -24\\ -3840} =0
\end{align}
\end{frame}
\begin{frame}
\frametitle{Solution}
Discriminant:
\begin{align}
D = (-24)^2 - 4 \times 3 \times (-3840) = 46656 \\
x = \frac{24 \pm \sqrt{46656}}{6} = \frac{24 \pm 216}{6}
\end{align}
Positive root:
\begin{align}
x = \frac{24 + 216}{6} = 40
\end{align}
Therefore,train speed = $40\frac{km}{h}$
\end{frame}

\begin{frame}[fragile]
\frametitle{C-Code}
\begin{lstlisting}[language=C]
// code.c
#include <stdio.h>

void get_roots(double *root1, double *root2) {
    double a = 3.0;
    double b = -24.0;
    double c = -3840.0;
    double D = b*b - 4*a*c;
    if (D < 0) {
        *root1 = 0.0;
        *root2 = 0.0;
    } else {
        double sqrtD = sqrt(D);
        *root1 = (-b + sqrtD) / (2*a);
        *root2 = (-b - sqrtD) / (2*a);
    }
}
\end{lstlisting}
\end{frame}
\begin{frame}[fragile]
\frametitle{Python-Code}
\begin{lstlisting}
# Code by GVV Sharma
# Modified for Problem Solution
# Released under GNU GPL
# Calculating area enclosed between curves
# call.py
import ctypes
import numpy as np

lib = ctypes.CDLL('./code.so')
\end{lstlisting}
\end{frame}
\begin{frame}[fragile]
\frametitle{Python-Code}
\begin{lstlisting}
get_roots = lib.get_roots
get_roots.argtypes = [
    ctypes.POINTER(ctypes.c_double),
    ctypes.POINTER(ctypes.c_double)
]
get_roots.restype = None

def get_quadratic_roots():
    root1 = ctypes.c_double()
    root2 = ctypes.c_double()
    get_roots(ctypes.byref(root1), ctypes.byref(root2))
    return root1.value, root2.value
\end{lstlisting}
\end{frame}

\end{document}